\subsection{Датчики в полёте и журнал}

Положение и ошибку мы уже читали (раздел~\ref{sec:python-measure}). Здесь --- остальные датчики Пионера и запись данных полёта.

\subsubsection*{Основные датчики}
\begin{itemize}
    \item \textbf{Барометр:} Измеряет высоту по давлению воздуха.
    \item \textbf{IMU (Гироскоп + Акселерометр):} Измеряет углы наклона и ускорения.
    \item \textbf{Дальномер (ToF):} Лазерный датчик расстояния (смотрит вниз), измеряет дистанцию до поверхности под дроном --- пола, крыши, холма. Команда \texttt{get\_dist\_sensor\_data()}, метры или \texttt{None}.
    \item \textbf{Система оптического потока (OptiFlow):} Камера, смотрящая вниз, для удержания позиции. В симуляторе её показания читает \texttt{get\_optical\_flow()} --- расширение симулятора: тройка (поток по X, поток по Y в рад/с в осях дрона, качество 0--255); скорость над поверхностью примерно равна потоку, умноженному на высоту.
\end{itemize}

\subsubsection*{Чтение данных дальномера}
Дальномер полезен для точного удержания высоты при посадке или полете над препятствиями.

\begin{codefile}[python]{python/get_telemetry/GET_dist_sensor_data.py}
\end{codefile}

Программа опрашивает дальномер бесконечно и печатает только настоящие измерения: \texttt{None} пропускается. Остановить её --- \key{Ctrl+C} или кнопкой остановки симулятора. Координаты относительно точки старта читает \texttt{get\_local\_position\_lps()}; над ровным полом её $z$ и показание дальномера почти совпадают, над холмом --- нет.

\begin{taskbox}[Мягкая посадка]
Взлетите на 2~м. Используя дальномер, снижайтесь командой \verb|set_manual_speed| с отрицательной \verb|vz| (раздел~\ref{sec:python-manual-speed}): 0.5~м/с, пока до поверхности больше 0.5~м, затем 0.1~м/с. Когда до поверхности меньше 0.15~м, отправьте нулевую скорость и \verb|land()|. Каждые 0.5~с печатайте показание дальномера; при \verb|None| скорость нулевая.
\expected{в журнале показание дальномера сначала падает быстро, после 0.5~м --- заметно медленнее; дрон касается площадки без удара, состояние становится \texttt{DISARMED}.}
\mission{Python}{65, 67}
\end{taskbox}

\begin{taskbox}[Черный ящик]
Только на компьютере (настоящий дрон или программа с компьютера на симулятор, раздел~\ref{sec:python-external}): в браузере файл программе недоступен. Во время полёта по квадрату каждые 0.5~с записывайте в файл \texttt{flight\_log.txt} строку «время;x;y;z;напряжение», а при \texttt{None} --- слово \texttt{нет} вместо числа. Файл открывайте конструкцией \verb|with open(...) as f:|, чтобы он закрылся и при ошибке.
\expected{после посадки в файле по строке на каждые полсекунды полёта, время растёт, координаты повторяют вершины квадрата; файл открывается в любом текстовом редакторе.}
\mission{Python}{72}
\end{taskbox}
